\begin{lemma}{Elementare Eigenschaften quadratintegrierbarer Abbildungen}
             {ElementareEigenschaftenQuadratintegrierbarerAbbildungen}

Seien $n \in \N$ und $\gamma \in L_2(\Omega, \mfa, \Prim, \R^n)$ zentriert,
unkorreliert und $L_2$--normiert.\\
Dann gilt:
\begin{enumerate}[(1)]
	\item \label{ElementareEigenschaftenQuadratintegrierbarerAbbildungen1}
 $\forall\, i, j \in \haken n : i\neq j\folgt\sprod{\gamma_i}{\gamma_j}=0$.
\item\label{ElementareEigenschaftenQuadratintegrierbarerAbbildungen2}
 $\bigg\|\Sum{i=1}n \gamma_i\bigg\|_2^2=\Sum{i=1}n \norm{\gamma_i}_2^2= n$.
\item\label{ElementareEigenschaftenQuadratintegrierbarerAbbildungen3}
 $\forall\, c \in \R^n: \norma{\bprod c \gamma_{n, L_2}}_2=\norm c_{2, n}$.
\item\label{ElementareEigenschaftenQuadratintegrierbarerAbbildungen4}
 $\norm\cdot_{2, n} \circ\gamma$ ist $\mfa-\mfb$--messbar und
    $\E(\norm\cdot_{2, n} \circ \gamma) \le \sqrt n$.
\end{enumerate}
\end{lemma}
\begin{proof}
(1) Seien $i \neq j \in \haken n$.\\
Dann ist $\sprod {\gamma_i} {\gamma_j}
= \E(\gamma_i\gamma_j) = \E(\gamma_i)\E(\gamma_j) = 0$, da $\gamma$ zentriert
ist.\\
(2) Die erste Gleichheit ist Pythagoras im Hilbertraum 
$L_2(\Omega, \mfa, \Prim, \R)$, die zweite folgt aus der Forderung, dass f�r
alle $i \in \haken n$ wegen $\gamma$ $L_2$--normiert $\norm{\gamma_i}_2 = 1$
gilt.\\
(3) Sei $c \in \R^n$. Dann gilt mit Pythagoras:
\[\norma{\bprod c \gamma _{n, L_2}}_2^2 = \norma{\sum_{i=1}^nc_i\gamma_i}_2^2
= \sum_{i=1}^n c_i^2 \norm{\gamma_i}_2^2 = \norm c _{2, n}^2\text.\]
(4) Sei $Y := \norm\cdot_{2,n}^2 \circ \gamma$.\\
Da $\gamma$ messbar,
$\norm\cdot_{2, n}$ stetig bez�glich sich selbst, $q: \R \rightarrow \R,
t\mapsto t^2$ stetig sind und auf $\R$ die Borel'sche $\sigma-$Algebra
vorliegt, sind $Y$ und $\sqrt\cdot \circ Y$ als Hintereinanderausf�hrung
messbarer und stetiger Funktionen selbst messbar.\\
Da $\norm\cdot_{2, n} \circ \gamma \ge 0$, gilt
$\norm\cdot_{2, n} \circ \gamma = \sqrt\cdot \circ Y$.\\
Die Wurzelabbildung ist konkav, also gilt mit der Jensenschen Ungleichung und
der Normiertheit der Komponentenfunktionen von $\gamma$:
\[\E\left(\sqrt\cdot \circ Y\right) \le \sqrt{\E(Y)} 
= \sqrt{\sum_{i=1}^n\E(\gamma_i^2)} = \sqrt n\text.\]
\end{proof}